\chapter{The seed}

The seed phrase is the wallet. Every key the device
derives comes from it. Anyone who sees it can empty
every account later, quietly, from another country.

If you even suspect it was seen, photographed, typed
into a connected device, or copied to a clipboard, treat
it as gone. Make a new wallet. Move the funds. Do not
``keep an eye on it.''

\section{Offline or it is not a backup}

Write it down when the device shows it. Paper is a
start. Metal survives fire and water better. Store
copies in places that do not burn, flood, or get
burglarized together.

Wallet vendors cannot recover it for you. That is the
point.

\section{What never happens}

No photo. No cloud. No chat. No email to yourself. No
plaintext file. No password manager. No ``I will
memorize it.'' Nobody legitimate will ask for it.
Conference USB sticks are not storage.

Do not travel with the full seed. Do not store every
piece in the same drawer.

\section{Splits are examples}

The source describes a 3-piece overlap split: words
1--16, 9--24, and 1--8 plus 17--24, so any two pieces
rebuild a 24-word seed. Shamir, vendor multi-share, and
a 25th word exist too.

Those are examples, not a policy. A split you cannot
rehearse is how you lock yourself out. A split stored
with people who do not know the code word is a story,
not a recovery plan. If you encrypt the backup, the
passphrase becomes another seed. Losing that is
permanent.

\section{Orgs and succession}

If this seed is a signer in a quorum, losing you should
not lose the quorum. Geographic copies. Theft of one
location must not rebuild the wallet. Operator loss must
not either.

Write down who can open which location, with a code
word, and test it. Next of kin needs a path that is not
a tweet. A lawyer's sealed instructions beat a puzzle
only you understood.

Revisit locations twice a year. Know where every copy
is. If the primary device died tonight, could you
recover before the rent is due?
